\EMhold

\EMsection{The One-Sided Laplace Transform}{The One-Sided Laplace Transform}

\begin{EMdefinition}{The one-sided Laplace transform}

\EMdisp{F(s)\EMbr =\mathcal{L}\{f(t)\}\triangleq\int_0^\infty f(t)\,e^{-st}\,dt\qquad\EMbr \EMbr \Longleftrightarrow\qquad\EMbr  f(t)\EMbr =\mathcal{L}^{-1}\{F(s)\}}

\end{EMdefinition}

The Laplace transform is one of the important tools in various engineering and mathematical applications. One of its applications is solving differential equations (read Chapter 6 of the course textbook).

The Laplace transform is a linear operator:
\EMdisp{\mathcal{L}\{af(t)+bg(t)\}\EMbr =a\,\mathcal{L}\{f(t)\}+b\,\mathcal{L}\{g(t)\}}

\begin{EMexample}{Example 1}

\EMdisp{f(t)\EMbr =\begin{cases}1&t\ge0\\ 0&t<0\end{cases}\ \EMbr \Longrightarrow\ \mathcal{L}\{f(t)\}\EMbr =F(s)\EMbr =\;?}

\EMdisp{\mathcal{L}\{f(t)\}\EMbr =\int_0^\infty1\times e^{-st}\,dt\EMbr =-\frac1se^{-st}\Big|_0^\infty\EMbr =\frac1s}

\end{EMexample}

\begin{EMexample}{Example 2}

\EMdisp{f(t)\EMbr =\begin{cases}e^{at}&t\ge0\\ 0&t<0\end{cases}\ \EMbr \Longrightarrow\ \mathcal{L}\{f(t)\}\EMbr =F(s)\EMbr =\;?}

\EMdisp{\mathcal{L}\{f(t)\}\EMbr =\int_0^\infty e^{at}\times e^{-st}\,dt\EMbr =\int_0^\infty e^{-(s-a)t}\,dt\EMbr =-\frac1{s-a}e^{-(s-a)t}\Big|_0^\infty\EMbr =\frac1{s-a}}

\end{EMexample}

\begin{EMexample}{Example 3}

\EMdisp{\mathcal{L}\{\cosh(at)\}\EMbr =\mathcal{L}\left\{\tfrac12e^{at}+\tfrac12e^{-at}\right\}\EMbr =\tfrac12\mathcal{L}\{e^{at}\}+\tfrac12\mathcal{L}\{e^{-at}\}\EMbr =\frac12\frac1{s-a}+\frac12\frac1{s+a}\EMbr =\frac s{s^2-a^2}}

\end{EMexample}

\begin{EMexample}{Example 4}

\EMdisp{\mathcal{L}\{\sinh(at)\}\EMbr =\mathcal{L}\left\{\tfrac12e^{at}-\tfrac12e^{-at}\right\}\EMbr =\tfrac12\mathcal{L}\{e^{at}\}-\tfrac12\mathcal{L}\{e^{-at}\}\EMbr =\frac12\frac1{s-a}-\frac12\frac1{s+a}\EMbr =\frac a{s^2-a^2}}

\end{EMexample}

\begin{EMexample}{Example 5}

Using \EMim{\displaystyle\int e^{ax}\cos nx\,dx=\frac{e^{ax}}{a^2+n^2}\big(a\cos nx+n\sin nx\big)}:
\EMdisp{\mathcal{L}\{\cos\omega t\}\EMbr =\int_0^\infty\cos\omega t\,e^{-st}\,dt\EMbr =\frac{e^{-st}}{s^2+\omega^2}\big(-s\cos\omega t+\omega\sin\omega t\big)\Big|_0^\infty\EMbr =\frac s{s^2+\omega^2}}

\end{EMexample}

\begin{EMexample}{Example 6}

Using \EMim{\displaystyle\int e^{ax}\sin nx\,dx=\frac{e^{ax}}{a^2+n^2}\big(a\sin nx-n\cos nx\big)}:
\EMdisp{\mathcal{L}\{\sin\omega t\}\EMbr =\int_0^\infty\sin\omega t\,e^{-st}\,dt\EMbr =\frac{e^{-st}}{s^2+\omega^2}\big(-s\sin\omega t-\omega\cos\omega t\big)\Big|_0^\infty\EMbr =\frac\omega{s^2+\omega^2}}

\end{EMexample}

\EMsubsection{Table of Laplace Transforms (Table 6.1 of the textbook; a longer table
is in Section 6.9)}{Table of Laplace Transforms (Table 6.1 of the textbook; a longer table is in Section 6.9)}

\Needspace{15\baselineskip}
\begin{longtable}[c]{@{}
  >{\centering\arraybackslash}p{(\columnwidth - 12\tabcolsep) * \real{0.1515}}
  >{\centering\arraybackslash}p{(\columnwidth - 12\tabcolsep) * \real{0.1515}}
  >{\centering\arraybackslash}p{(\columnwidth - 12\tabcolsep) * \real{0.1515}}
  >{\raggedright\arraybackslash}p{(\columnwidth - 12\tabcolsep) * \real{0.0909}}
  >{\centering\arraybackslash}p{(\columnwidth - 12\tabcolsep) * \real{0.1515}}
  >{\centering\arraybackslash}p{(\columnwidth - 12\tabcolsep) * \real{0.1515}}
  >{\centering\arraybackslash}p{(\columnwidth - 12\tabcolsep) * \real{0.1515}}@{}}
\toprule\noalign{}
\rowcolor{EMindigoL}\begin{minipage}[b]{\linewidth}\centering
\end{minipage} & \begin{minipage}[b]{\linewidth}\centering
\EMim{f(t)}
\end{minipage} & \begin{minipage}[b]{\linewidth}\centering
\EMim{\mathcal{L}(f)}
\end{minipage} & \begin{minipage}[b]{\linewidth}\raggedright
\end{minipage} & \begin{minipage}[b]{\linewidth}\centering
\end{minipage} & \begin{minipage}[b]{\linewidth}\centering
\EMim{f(t)}
\end{minipage} & \begin{minipage}[b]{\linewidth}\centering
\EMim{\mathcal{L}(f)}
\end{minipage} \\
\midrule\noalign{}
\endhead
\bottomrule\noalign{}
\endlastfoot
1 & \EMim{1} & \EMim{1/s} & & 7 & \EMim{\cos\omega t} & \EMim{\dfrac{s}{s^2+\omega^2}} \\
2 & \EMim{t} & \EMim{1/s^2} & & 8 & \EMim{\sin\omega t} & \EMim{\dfrac{\omega}{s^2+\omega^2}} \\
3 & \EMim{t^2} & \EMim{2!/s^3} & & 9 & \EMim{\cosh at} & \EMim{\dfrac{s}{s^2-a^2}} \\
4 & \EMim{t^n} \EMim{(n=0,1,\dots)} & \EMim{\dfrac{n!}{s^{n+1}}} & & 10 & \EMim{\sinh at} & \EMim{\dfrac{a}{s^2-a^2}} \\
5 & \EMim{t^a} (\EMim{a} positive) & \EMim{\dfrac{\Gamma(a+1)}{s^{a+1}}} & & 11 & \EMim{e^{at}\cos\omega t} & \EMim{\dfrac{s-a}{(s-a)^2+\omega^2}} \\
6 & \EMim{e^{at}} & \EMim{\dfrac1{s-a}} & & 12 & \EMim{e^{at}\sin\omega t} & \EMim{\dfrac{\omega}{(s-a)^2+\omega^2}} \\
\end{longtable}

\begin{EMunit}

\EMdisp{\Gamma(x)\EMbr =\int_0^\infty e^{-t}t^{x-1}\,dt}

\end{EMunit}

\EMsubsection{Properties of the One-Sided Laplace Transform (the table of Section 6.8
of the textbook)}{Properties of the One-Sided Laplace Transform (the table of Section 6.8 of the textbook)}

With the unit step function
\EMdisp{u(t)\triangleq\begin{cases}1&t\ge0\\ 0&t<0\end{cases}}

\Needspace{22\baselineskip}
\begin{longtable}[c]{@{}
  >{\raggedright\arraybackslash}p{(\columnwidth - 2\tabcolsep) * \real{0.4444}}
  >{\centering\arraybackslash}p{(\columnwidth - 2\tabcolsep) * \real{0.5556}}@{}}
\toprule\noalign{}
\rowcolor{EMindigoL}\begin{minipage}[b]{\linewidth}\raggedright
Property
\end{minipage} & \begin{minipage}[b]{\linewidth}\centering
Formula
\end{minipage} \\
\midrule\noalign{}
\endhead
\bottomrule\noalign{}
\endlastfoot
Definition of the Laplace transform & \EMim{F(s)=\mathcal{L}\{f(t)\}\triangleq\displaystyle\int_0^\infty f(t)e^{-st}\,dt} \\
Inverse Laplace transform & \EMim{f(t)=\mathcal{L}^{-1}\{F(s)\}} \\
Linearity & \EMim{\mathcal{L}\{af(t)+bg(t)\}=a\mathcal{L}\{f(t)\}+b\mathcal{L}\{g(t)\}} \\
Time shift & \EMim{\mathcal{L}\{f(t-a)u(t-a)\}=e^{-as}F(s)} \\
Shift in the \EMim{s} domain & \EMim{\mathcal{L}\{e^{at}f(t)\}=F(s-a)} \\
Time scaling & \EMim{\mathcal{L}\{f(at)\}=\dfrac1aF(s/a)} \\
Time derivative & \EMim{\mathcal{L}\{f^{(n)}(t)\}=s^nF(s)-s^{n-1}f(0)-s^{n-2}f^{(1)}(0)-\cdots-f^{(n-1)}(0)} \\
Time integral & \EMim{\mathcal{L}\left\{\displaystyle\int_0^tf(\tau)\,d\tau\right\}=\dfrac1sF(s)} \\
Derivative in the \EMim{s} domain & \EMim{\mathcal{L}\{tf(t)\}=-F'(s)} \\
Integral in the \EMim{s} domain & \EMim{\mathcal{L}\left\{\dfrac{f(t)}t\right\}=\displaystyle\int_s^\infty F(\nu)\,d\nu} \\
\end{longtable}

\EMhold

\EMsection{Using the Laplace Transform to Solve Differential Equations}{Using the Laplace Transform to Solve Differential Equations}

\begin{EMunit}

The steps for solving a differential equation with initial conditions using the Laplace transform:

\begin{enumerate}
\def\labelenumi{\arabic{enumi}.}
\tightlist
\item
  convert the initial-value problem into an algebraic problem;
\item
  compute the solution of the algebraic problem;
\item
  solve the initial-value problem (return the solution to the time domain).
\end{enumerate}

\EMfigure{m07-solve-scheme}{The steps of solving an initial-value problem with the Laplace
transform: initial-value problem \EMim{\xrightarrow{1}} algebraic
problem \EMim{\xrightarrow{2}} solution of the algebraic problem
\EMim{\xrightarrow{3}} solution of the initial-value problem}

\end{EMunit}

\EMhold

\EMsubsection{Ordinary Differential Equations}{Ordinary Differential Equations}

\begin{EMexample}{}

Solve the differential equation \EMim{y''-y=t}, \EMim{y(0)=1}, \EMim{y'(0)=1} using the Laplace transform.

\EMdisp{s^2Y-sy(0)-y'(0)-Y\EMbr =\frac1{s^2}\ \EMbr \Longrightarrow\ (s^2-1)Y(s)\EMbr =s+1+\frac1{s^2}}
\EMdisp{Y(s)\EMbr =\frac{s+1}{s^2-1}+\frac1{s^2(s^2-1)}\EMbr =\frac1{s-1}+\left(\frac1{s^2-1}-\frac1{s^2}\right)}
\EMdisp{y(t)\EMbr =\mathcal{L}^{-1}\{Y(s)\}\EMbr =\mathcal{L}^{-1}\left\{\frac1{s-1}\right\}+\mathcal{L}^{-1}\left\{\frac1{s^2-1}\right\}-\mathcal{L}^{-1}\left\{\frac1{s^2}\right\}\EMbr =e^tu(t)+\sinh t\,u(t)-t\,u(t)}

\end{EMexample}

\EMhold

\EMsubsection{Partial Differential Equations}{Partial Differential Equations}

\begin{EMexample}{}

Solve the following partial differential equation using the Laplace transform.
\EMdisp{\frac{\partial w}{\partial x}+x\frac{\partial w}{\partial t}\EMbr =0,\qquad\EMbr  w(x,0)\EMbr =0,\qquad\EMbr  w(0,t)\EMbr =t\,u(t)}

We take the Laplace transform of both sides with respect to \EMim{t}:
\EMdisp{\mathcal{L}_t\left\{\frac{\partial w}{\partial x}+x\frac{\partial w}{\partial t}\right\}\EMbr =\mathcal{L}_t\left\{\frac{\partial w}{\partial x}\right\}+x\,\mathcal{L}_t\left\{\frac{\partial w}{\partial t}\right\}\EMbr =0}
\EMdisp{\frac\partial{\partial x}w(x,s)+x\big[s\,w(x,s)-w(x,0)\big]\EMbr =0,\qquad\EMbr  w(x,0)\EMbr =0}
\EMdisp{\frac\partial{\partial x}w(x,s)+xs\,w(x,s)\EMbr =0\ \EMbr \Rightarrow\ \frac{dw}w\EMbr =-sx\,dx\ \EMbr \Rightarrow\ \ln w\EMbr =-s\frac{x^2}2+\ln C\ \EMbr \Rightarrow\ w(x,s)\EMbr =Ce^{-\frac{x^2}2s}}
From \EMim{w(0,t)=t\,u(t)} we have \EMim{w(0,s)=\mathcal{L}\{w(0,t)\}=1/s^2=C}:
\EMdisp{w(x,s)\EMbr =\frac1{s^2}e^{-\frac{x^2}2s}}
Using the time-shift property \EMim{\mathcal{L}\{f(t-a)u(t-a)\}=e^{-as}F(s)}:
\EMdisp{w(x,t)\EMbr =\left(t-\frac{x^2}2\right)\times u\!\left(t-\frac{x^2}2\right)}

\end{EMexample}
